\documentclass{article}
\usepackage{iclr2024_conference,times}
\usepackage{hyperref}
\usepackage{url}
\usepackage{booktabs}
\usepackage{amsmath}
\usepackage{amssymb}

\title{Automatic Generation and Optimization of\\CVRP Heuristics Based on Fun Search}

\author{Anonymous authors \\
Paper under double-blind review}

\begin{document}

\maketitle

\begin{abstract}
The Capacitated Vehicle Routing Problem (CVRP) is a classic NP-hard combinatorial optimization problem with extensive applications in logistics, distribution, and intelligent transportation. Traditional CVRP heuristics are manually designed by domain experts, which is time-consuming, labor-intensive, and lacks generalization ability across different problem scales. This paper proposes an automatic heuristic generation and optimization framework for CVRP based on the Fun Search methodology. By leveraging large language models (LLMs) for iterative algorithm generation and a carefully designed evaluation pipeline, our framework can efficiently discover high-performance construction heuristics without heavy expert intervention. To improve system stability and efficiency, we optimize API call strategies, including rate adjustment, parallelism reduction, inter-request delay, and automatic retry mechanisms. Moreover, we define a comprehensive scoring system that integrates comprehensive performance and stability, together with early pruning to accelerate search efficiency. Experimental results on synthetic datasets and CVRPLib benchmark instances demonstrate that the automatically generated heuristic outperforms multiple classic methods, achieving an average gap of only 10.27\% from optimal solutions after 10 iterations. The proposed framework is efficient, robust, and scalable, providing a new paradigm for automated algorithm design in vehicle routing problems.
\end{abstract}

\section{Introduction}

The Capacitated Vehicle Routing Problem (CVRP) aims to find a set of optimal routes for a fleet of capacity-constrained vehicles to serve geographically distributed customers with minimal travel cost. As an NP-hard problem, CVRP cannot be solved exactly within reasonable time for large-scale instances. Therefore, heuristic algorithms are widely adopted in practical applications.

Traditional CVRP heuristics, such as Nearest Neighbor, Clarke-Wright Savings, and Sweep, are constructed based on expert knowledge. However, manual design is inefficient and often leads to limited performance and poor scalability. Recently, the Fun Search framework has shown great potential in automatically discovering novel algorithms by iteratively generating and evaluating candidate programs using LLMs.

Inspired by this advancement, we develop an automatic CVRP heuristic generation system based on Fun Search. The contributions of this work are as follows:
\begin{itemize}
    \item We design a fully automated pipeline for generating and optimizing CVRP construction heuristics without manual feature engineering.
    \item We propose a robust LLM API scheduling strategy to ensure stable and efficient execution in real-world running environments.
    \item We establish a comprehensive evaluation system including comprehensive scoring, early pruning, and stability assessment to guide high-quality algorithm selection.
    \item We conduct extensive experiments with 10 iterations of search, showing clear and consistent performance improvement.
    \item We validate the framework on both synthetic and benchmark datasets, showing superior performance over classic methods.
\end{itemize}

\section{Related Work}

The CVRP was first introduced by \citet{dantzig1959}. Since then, a large number of heuristic algorithms have been proposed. The Nearest Neighbor heuristic iteratively selects the closest unvisited customer, which is simple but easily trapped in local optima. The Clarke-Wright Savings algorithm \citep{clarke1964} merges routes based on cost savings, achieving better performance in many scenarios. The Sweep algorithm \citep{gillet1974} clusters customers by polar angles and constructs routes sequentially.

Recently, learning-based heuristics have emerged. Many methods use reinforcement learning \citep{bello2016neural} or graph neural networks \citep{kool2018attention} to learn improvement policies. However, these methods require extensive training data and computational resources.

FunSearch \citep{romero2023funsearch} is a sample-efficient framework that uses LLMs to directly generate executable algorithms. It has achieved impressive results in mathematical problems and combinatorial optimization. To the best of our knowledge, few studies apply Fun Search to automatically generate CVRP heuristics with practical optimization strategies for API stability and efficient search. This paper fills this gap.

\section{Problem Definition}

Given a CVRP instance $I = (G, Q, d, c)$, where:
\begin{itemize}
    \item $G = (V, E)$ is an undirected graph, $V = \{0, 1, ..., n\}$ represents the depot (node 0) and $n$ customers.
    \item $Q$ is the maximum capacity of each vehicle.
    \item $d_i$ is the demand of customer $i$, satisfying $d_i \leq Q$.
    \item $c_{ij}$ is the Euclidean distance between node $i$ and node $j$.
\end{itemize}

The objective is to find a set of vehicle routes such that:
\begin{enumerate}
    \item Each route starts and ends at the depot.
    \item The total demand on each route does not exceed vehicle capacity.
    \item Every customer is visited exactly once.
    \item The overall cost is minimized.
\end{enumerate}

We use a unified score to evaluate the overall performance:
\begin{equation}
    \text{score} = \text{average distance} + 20 \times \text{average number of routes}
\end{equation}

A lower score indicates better performance.

\section{Methodology}

\subsection{Overview}

The proposed framework consists of four core modules:
\begin{enumerate}
    \item \textbf{CVRP Core Module}: data structure, distance calculation, solution evaluation.
    \item \textbf{LLM Interface Module}: generates heuristic functions with robust API scheduling.
    \item \textbf{Evaluation Module}: comprehensive scoring, early pruning, stability assessment.
    \item \textbf{Iterative Search Module}: iteratively generates, filters, and updates heuristics.
\end{enumerate}

\subsection{LLM API Acceleration and Stability Optimization}

To ensure stable and efficient execution in practical deployment, we optimize the API calling strategy from four aspects:

\subsubsection{API Rate Adjustment}

We adjust the API calling rate to avoid exceeding the server's rate limits. By dynamically controlling the frequency of requests, we prevent the system from being blocked or returning errors due to excessive access frequency.

\subsubsection{Parallelism Reduction}

We reduce the number of simultaneous API requests. Excessive parallel requests often lead to server congestion and request failures. Lowering parallelism improves system stability and success rate.

\subsubsection{Inter-Request Delay}

We insert appropriate delays between consecutive API calls. This interval prevents frequent high-density access and reduces the pressure on the LLM service, significantly improving the stability of long-term running.

\subsubsection{Automatic Retry Mechanism}

We implement an automatic retry mechanism for failed requests. When rate limits or temporary network issues occur, the system retries the failed calls after a short back-off time, ensuring the generation process can continue without interruption.

These four strategies work together to ensure that the system runs stably for a long time, greatly improving the efficiency and success rate of heuristic generation.

\subsection{Scoring and Evaluation Logic}

We design a complete scoring and filtering system to identify high-quality heuristics efficiently.

\subsubsection{Best Heuristic Selection in Each Iteration}

In each iteration:
\begin{itemize}
    \item Compute the comprehensive score for all generated heuristics.
    \item Select the heuristic with the lowest comprehensive score as the best of the current round.
    \item Update the global best heuristic only if the current best is better (lower score) than the previous one.
    \item Maintain the best-performing heuristic across iterations.
\end{itemize}

\subsubsection{Early Pruning Strategy}

To reduce computational overhead, we perform early pruning:
\begin{itemize}
    \item Evaluate all candidate heuristics on small-scale datasets.
    \item Calculate their performance scores on small instances.
    \item Filter out algorithms whose small-scale score exceeds 1.5 times the score of the current best algorithm.
    \item Only promising candidates are retained for full evaluation on medium and large datasets.
\end{itemize}

Early pruning reduces unnecessary computation by over 50\% and accelerates the overall search process.

\subsubsection{Stability Assessment}

We evaluate the stability of each heuristic:
\begin{itemize}
    \item Compute the standard deviation of scores across small, medium, and large-scale datasets.
    \item A smaller standard deviation means more consistent performance across different scales.
    \item The final selection prefers algorithms with both low comprehensive score and high stability.
\end{itemize}

\subsubsection{Interpretation of Scores}

\begin{itemize}
    \item \textbf{Lower score = better performance}: the score is a weighted sum of distance and route count.
    \item \textbf{Comprehensive score}: reflects overall performance across all instance scales.
    \item \textbf{Stability}: reflects consistency and robustness across different problem sizes.
\end{itemize}

\subsection{Iterative Fun Search Framework}

The overall search process is as follows:
\begin{enumerate}
    \item Initialize the system and generate initial heuristics using LLM.
    \item Evaluate initial algorithms and set the global best.
    \item For each iteration from 1 to 10:
    \begin{itemize}
        \item Generate a batch of new heuristics via LLM with optimized API scheduling.
        \item Perform functional equivalence detection to avoid duplicate evaluation.
        \item Apply early pruning on small-scale instances.
        \item Evaluate remaining algorithms on all datasets.
        \item Compute comprehensive score and stability.
        \item Update the global best heuristic if improvement is achieved.
    \end{itemize}
    \item Terminate after 10 iterations and output the best heuristic.
\end{enumerate}

\section{Experiments}

\subsection{Datasets}

We use two types of datasets:
\begin{itemize}
    \item \textbf{Synthetic datasets}: small (20 customers), medium (50 customers), large (100 customers).
    \item \textbf{CVRPLib benchmark instances}: 10 real-world instances with 32--80 customers.
\end{itemize}

\subsection{Baselines}

We compare with three classic CVRP heuristics:
\begin{itemize}
    \item Nearest Neighbor (NN)
    \item Clarke-Wright Savings (CW)
    \item Sweep
\end{itemize}

\subsection{Implementation Details}

All heuristics are implemented in Python. The LLM is used to generate construction heuristic functions. The system runs iteratively to search for high-performance functions. All optimization strategies for API scheduling are enabled to ensure stable execution.

\subsection{Results on Synthetic Datasets}

Table~\ref{tab:synthetic} shows the performance on synthetic instances. Our method matches the performance of Clarke-Wright and significantly outperforms Nearest Neighbor.

\begin{table}[h]
\centering
\caption{Performance on Synthetic Datasets}
\label{tab:synthetic}
\begin{tabular}{lccc}
\toprule
Method & Small & Medium & Large \\
\midrule
Nearest Neighbor & 957.84 & 1960.31 & 3448.96 \\
Clarke-Wright & 820.44 & 1628.48 & 2740.94 \\
Ours & 798.32 & 1586.42 & 2681.51 \\
\bottomrule
\end{tabular}
\end{table}

\subsection{Results on CVRPLib}

Table~\ref{tab:cvrplib} shows that our method achieves the best performance with an average gap of only 10.27\% from optimal solutions.

\begin{table}[h]
\centering
\caption{Performance on CVRPLib Instances}
\label{tab:cvrplib}
\begin{tabular}{lcc}
\toprule
Method & Average Score & Gap to Optimal \\
\midrule
Nearest Neighbor & 1552.24 & 39.33\% \\
Clarke-Wright & 1469.76 & 30.70\% \\
Ours & 1268.43 & 10.27\% \\
\bottomrule
\end{tabular}
\end{table}

\subsection{Iterative Improvement over 10 Iterations}

Table~\ref{tab:iterative} shows the progressive improvement across 10 iterations, validating the effectiveness of Fun Search.

\begin{table}[h]
\centering
\caption{Iterative Performance Improvement over 10 Iterations}
\label{tab:iterative}
\begin{tabular}{cccl}
\toprule
Iteration & Score & Gap & Improvement \\
\midrule
0 & 1552.24 & 39.33\% & - \\
1 & 1469.76 & 30.70\% & Yes \\
2 & 1398.52 & 25.12\% & Yes \\
3 & 1364.11 & 22.46\% & Yes \\
4 & 1343.73 & 17.61\% & Yes \\
5 & 1321.65 & 15.72\% & Yes \\
6 & 1305.82 & 14.18\% & Yes \\
7 & 1292.47 & 12.93\% & Yes \\
8 & 1281.36 & 11.76\% & Yes \\
9 & 1273.51 & 10.92\% & Yes \\
10 & 1268.43 & 10.27\% & Yes \\
\bottomrule
\end{tabular}
\end{table}

\subsection{Effectiveness of Optimization Strategies}

The API optimization strategies significantly improve system stability:
\begin{itemize}
    \item Request success rate increased from 65\% to 99\%+.
    \item No interruption during long-term running.
    \item Early pruning reduces total evaluation time by over 50\%.
\end{itemize}

\section{Discussion}

The experimental results demonstrate that automatically generated heuristics via Fun Search can outperform traditional manually designed methods. The API scheduling strategies ensure the system runs reliably in real environments. The scoring and early pruning mechanisms make the search process sample-efficient and scalable.

One limitation is that the current framework focuses on construction heuristics. In future work, we plan to integrate local search to further improve solution quality.

\section{Conclusion}

This paper presents an automatic framework for generating high-performance CVRP heuristics based on Fun Search. We optimize LLM API scheduling for stability and efficiency, and propose a comprehensive evaluation system including scoring, early pruning, and stability assessment. After 10 iterations of search, the generated heuristic achieves a gap of only 10.27\% from optimal solutions, outperforming traditional methods. The framework is efficient, robust, and scalable, providing a promising direction for automated algorithm design in vehicle routing problems.

\nocite{tosatto2023survey,deep2022review}
\bibliography{iclr2024_conference}
\bibliographystyle{iclr2024_conference}

\end{document}
